\documentclass[12pt,letterpaper]{article}
\usepackage[margin=0in]{geometry}
\usepackage{tikz}
\usepackage[T1]{fontenc}
\usepackage{lmodern}\pdfmapfile{+lm.map}
\pagestyle{empty}
\setlength{\parindent}{0pt}\hyphenpenalty=10000\exhyphenpenalty=10000
\begin{document}
\null
\begin{tikzpicture}[remember picture,overlay,x=1mm,y=-1mm]
\begin{scope}[shift={(current page.north west)}]
\node[anchor=north west,inner sep=0pt,outer sep=0pt,align=left,text width=187mm] at (15,12) {\fontsize{11}{13.86}\selectfont Week 26 / Same area, different boundaries / Grades 2--3};
\node[anchor=north west,inner sep=0pt,outer sep=0pt,align=left,text width=184mm] at (15,26) {\fontsize{12}{15.12}\selectfont Keep all tiles flat and in one piece, joined along whole sides without overlaps. Count every side with no tile beside it, including sides around holes. Turns and flips count as the same shape.};
\node[anchor=north west,inner sep=0pt,outer sep=0pt,align=left,text width=184mm] at (15,66) {\fontsize{14}{17.64}\selectfont \textbf{Problem 1:} Use eight tiles to make a shape with the shortest boundary you can and one with the longest boundary you can. Find a different shape reaching each of those lengths. Draw all four shapes and record their boundary lengths.};
\begin{scope}[shift={(15,107)}]
\draw[step=7.5mm,black!27,line width=.3pt] (0,0) grid (75.0,45.0);
\end{scope}
\draw[black!27,line width=.3pt] (15,107) rectangle (90.0,152.0);
\node[anchor=north west,inner sep=0pt,outer sep=0pt,align=left] at (15,155) {\fontsize{12}{15.12}\selectfont Shortest boundary: \rule{12mm}{.3pt}};
\begin{scope}[shift={(114,107)}]
\draw[step=7.5mm,black!27,line width=.3pt] (0,0) grid (75.0,45.0);
\end{scope}
\draw[black!27,line width=.3pt] (114,107) rectangle (189.0,152.0);
\node[anchor=north west,inner sep=0pt,outer sep=0pt,align=left] at (114,155) {\fontsize{12}{15.12}\selectfont Longest boundary: \rule{12mm}{.3pt}};
\begin{scope}[shift={(15,184)}]
\draw[step=7.5mm,black!27,line width=.3pt] (0,0) grid (75.0,45.0);
\end{scope}
\draw[black!27,line width=.3pt] (15,184) rectangle (90.0,229.0);
\node[anchor=north west,inner sep=0pt,outer sep=0pt,align=left] at (15,232) {\fontsize{12}{15.12}\selectfont Shortest boundary: \rule{12mm}{.3pt}};
\begin{scope}[shift={(114,184)}]
\draw[step=7.5mm,black!27,line width=.3pt] (0,0) grid (75.0,45.0);
\end{scope}
\draw[black!27,line width=.3pt] (114,184) rectangle (189.0,229.0);
\node[anchor=north west,inner sep=0pt,outer sep=0pt,align=left] at (114,232) {\fontsize{12}{15.12}\selectfont Longest boundary: \rule{12mm}{.3pt}};
\node[anchor=north west,inner sep=0pt,outer sep=0pt,align=left] at (15,268) {\fontsize{9}{11.34}\selectfont Bellingham Math Circle / Week 26 / W26-23-v2};
\node[anchor=north east,inner sep=0pt,outer sep=0pt,align=left] at (199,268) {\fontsize{9}{11.34}\selectfont 1};
\end{scope}
\end{tikzpicture}
\newpage
\null
\begin{tikzpicture}[remember picture,overlay,x=1mm,y=-1mm]
\begin{scope}[shift={(current page.north west)}]
\node[anchor=north west,inner sep=0pt,outer sep=0pt,align=left,text width=187mm] at (15,12) {\fontsize{11}{13.86}\selectfont Week 26 / Same area, different boundaries / Grades 2--3};
\node[anchor=north west,inner sep=0pt,outer sep=0pt,align=left,text width=184mm] at (15,28) {\fontsize{14}{17.64}\selectfont \textbf{Problem 2:} Add one tile to each shape. Find every different change you can make to its boundary length.};
\filldraw[fill=black!13,draw=black,line width=.75pt] (15,65) rectangle ++(8,8);
\filldraw[fill=black!13,draw=black,line width=.75pt] (23,65) rectangle ++(8,8);
\filldraw[fill=black!13,draw=black,line width=.75pt] (31,65) rectangle ++(8,8);
\filldraw[fill=black!13,draw=black,line width=.75pt] (39,65) rectangle ++(8,8);
\filldraw[fill=black!13,draw=black,line width=.75pt] (47,65) rectangle ++(8,8);
\node[anchor=north west,inner sep=0pt,outer sep=0pt,align=left] at (15,95) {\fontsize{11}{13.86}\selectfont Before: \rule{12mm}{.3pt}};
\begin{scope}[shift={(105,63)}]
\draw[step=6.5mm,black!27,line width=.3pt] (0,0) grid (39.0,26.0);
\end{scope}
\draw[black!27,line width=.3pt] (105,63) rectangle (144.0,89.0);
\node[anchor=north west,inner sep=0pt,outer sep=0pt,align=left] at (105,94) {\fontsize{11}{13.86}\selectfont After: \rule{12mm}{.3pt}};
\node[anchor=north west,inner sep=0pt,outer sep=0pt,align=left] at (105,100) {\fontsize{11}{13.86}\selectfont Change: \rule{10mm}{.3pt}};
\begin{scope}[shift={(153,63)}]
\draw[step=6.5mm,black!27,line width=.3pt] (0,0) grid (39.0,26.0);
\end{scope}
\draw[black!27,line width=.3pt] (153,63) rectangle (192.0,89.0);
\node[anchor=north west,inner sep=0pt,outer sep=0pt,align=left] at (153,94) {\fontsize{11}{13.86}\selectfont After: \rule{12mm}{.3pt}};
\node[anchor=north west,inner sep=0pt,outer sep=0pt,align=left] at (153,100) {\fontsize{11}{13.86}\selectfont Change: \rule{10mm}{.3pt}};
\filldraw[fill=black!13,draw=black,line width=.75pt] (15,114) rectangle ++(8,8);
\filldraw[fill=black!13,draw=black,line width=.75pt] (23,114) rectangle ++(8,8);
\filldraw[fill=black!13,draw=black,line width=.75pt] (23,122) rectangle ++(8,8);
\filldraw[fill=black!13,draw=black,line width=.75pt] (31,122) rectangle ++(8,8);
\filldraw[fill=black!13,draw=black,line width=.75pt] (31,130) rectangle ++(8,8);
\node[anchor=north west,inner sep=0pt,outer sep=0pt,align=left] at (15,144) {\fontsize{11}{13.86}\selectfont Before: \rule{12mm}{.3pt}};
\begin{scope}[shift={(105,112)}]
\draw[step=6.5mm,black!27,line width=.3pt] (0,0) grid (39.0,26.0);
\end{scope}
\draw[black!27,line width=.3pt] (105,112) rectangle (144.0,138.0);
\node[anchor=north west,inner sep=0pt,outer sep=0pt,align=left] at (105,143) {\fontsize{11}{13.86}\selectfont After: \rule{12mm}{.3pt}};
\node[anchor=north west,inner sep=0pt,outer sep=0pt,align=left] at (105,149) {\fontsize{11}{13.86}\selectfont Change: \rule{10mm}{.3pt}};
\begin{scope}[shift={(153,112)}]
\draw[step=6.5mm,black!27,line width=.3pt] (0,0) grid (39.0,26.0);
\end{scope}
\draw[black!27,line width=.3pt] (153,112) rectangle (192.0,138.0);
\node[anchor=north west,inner sep=0pt,outer sep=0pt,align=left] at (153,143) {\fontsize{11}{13.86}\selectfont After: \rule{12mm}{.3pt}};
\node[anchor=north west,inner sep=0pt,outer sep=0pt,align=left] at (153,149) {\fontsize{11}{13.86}\selectfont Change: \rule{10mm}{.3pt}};
\filldraw[fill=black!13,draw=black,line width=.75pt] (15,163) rectangle ++(8,8);
\filldraw[fill=black!13,draw=black,line width=.75pt] (23,163) rectangle ++(8,8);
\filldraw[fill=black!13,draw=black,line width=.75pt] (31,163) rectangle ++(8,8);
\filldraw[fill=black!13,draw=black,line width=.75pt] (15,171) rectangle ++(8,8);
\filldraw[fill=black!13,draw=black,line width=.75pt] (31,171) rectangle ++(8,8);
\node[anchor=north west,inner sep=0pt,outer sep=0pt,align=left] at (15,193) {\fontsize{11}{13.86}\selectfont Before: \rule{12mm}{.3pt}};
\begin{scope}[shift={(105,161)}]
\draw[step=6.5mm,black!27,line width=.3pt] (0,0) grid (39.0,26.0);
\end{scope}
\draw[black!27,line width=.3pt] (105,161) rectangle (144.0,187.0);
\node[anchor=north west,inner sep=0pt,outer sep=0pt,align=left] at (105,192) {\fontsize{11}{13.86}\selectfont After: \rule{12mm}{.3pt}};
\node[anchor=north west,inner sep=0pt,outer sep=0pt,align=left] at (105,198) {\fontsize{11}{13.86}\selectfont Change: \rule{10mm}{.3pt}};
\begin{scope}[shift={(153,161)}]
\draw[step=6.5mm,black!27,line width=.3pt] (0,0) grid (39.0,26.0);
\end{scope}
\draw[black!27,line width=.3pt] (153,161) rectangle (192.0,187.0);
\node[anchor=north west,inner sep=0pt,outer sep=0pt,align=left] at (153,192) {\fontsize{11}{13.86}\selectfont After: \rule{12mm}{.3pt}};
\node[anchor=north west,inner sep=0pt,outer sep=0pt,align=left] at (153,198) {\fontsize{11}{13.86}\selectfont Change: \rule{10mm}{.3pt}};
\filldraw[fill=black!13,draw=black,line width=.75pt] (15,212) rectangle ++(8,8);
\filldraw[fill=black!13,draw=black,line width=.75pt] (15,220) rectangle ++(8,8);
\filldraw[fill=black!13,draw=black,line width=.75pt] (15,228) rectangle ++(8,8);
\filldraw[fill=black!13,draw=black,line width=.75pt] (23,212) rectangle ++(8,8);
\filldraw[fill=black!13,draw=black,line width=.75pt] (23,228) rectangle ++(8,8);
\filldraw[fill=black!13,draw=black,line width=.75pt] (31,212) rectangle ++(8,8);
\filldraw[fill=black!13,draw=black,line width=.75pt] (31,220) rectangle ++(8,8);
\filldraw[fill=black!13,draw=black,line width=.75pt] (31,228) rectangle ++(8,8);
\node[anchor=north west,inner sep=0pt,outer sep=0pt,align=left] at (15,242) {\fontsize{11}{13.86}\selectfont Before: \rule{12mm}{.3pt}};
\begin{scope}[shift={(105,210)}]
\draw[step=6.5mm,black!27,line width=.3pt] (0,0) grid (39.0,26.0);
\end{scope}
\draw[black!27,line width=.3pt] (105,210) rectangle (144.0,236.0);
\node[anchor=north west,inner sep=0pt,outer sep=0pt,align=left] at (105,241) {\fontsize{11}{13.86}\selectfont After: \rule{12mm}{.3pt}};
\node[anchor=north west,inner sep=0pt,outer sep=0pt,align=left] at (105,247) {\fontsize{11}{13.86}\selectfont Change: \rule{10mm}{.3pt}};
\begin{scope}[shift={(153,210)}]
\draw[step=6.5mm,black!27,line width=.3pt] (0,0) grid (39.0,26.0);
\end{scope}
\draw[black!27,line width=.3pt] (153,210) rectangle (192.0,236.0);
\node[anchor=north west,inner sep=0pt,outer sep=0pt,align=left] at (153,241) {\fontsize{11}{13.86}\selectfont After: \rule{12mm}{.3pt}};
\node[anchor=north west,inner sep=0pt,outer sep=0pt,align=left] at (153,247) {\fontsize{11}{13.86}\selectfont Change: \rule{10mm}{.3pt}};
\node[anchor=north west,inner sep=0pt,outer sep=0pt,align=left] at (15,268) {\fontsize{9}{11.34}\selectfont Bellingham Math Circle / Week 26 / W26-23-v2};
\node[anchor=north east,inner sep=0pt,outer sep=0pt,align=left] at (199,268) {\fontsize{9}{11.34}\selectfont 2};
\end{scope}
\end{tikzpicture}
\newpage
\null
\begin{tikzpicture}[remember picture,overlay,x=1mm,y=-1mm]
\begin{scope}[shift={(current page.north west)}]
\node[anchor=north west,inner sep=0pt,outer sep=0pt,align=left,text width=187mm] at (15,12) {\fontsize{11}{13.86}\selectfont Week 26 / Same area, different boundaries / Grades 2--3};
\node[anchor=north west,inner sep=0pt,outer sep=0pt,align=left,text width=184mm] at (15,28) {\fontsize{14}{17.64}\selectfont \textbf{Problem 3:} A shared side is where two tiles meet. Use ten tiles to make shapes with 9, 10, 11, 12, and 13 shared sides. Record each boundary length. Find a rule that gives the boundary length from the number of tiles and the number of shared sides.};
\begin{scope}[shift={(15,77)}]
\draw[step=6.2mm,black!27,line width=.3pt] (0,0) grid (74.4,31.0);
\end{scope}
\draw[black!27,line width=.3pt] (15,77) rectangle (89.4,108.0);
\node[anchor=north west,inner sep=0pt,outer sep=0pt,align=left] at (15,111) {\fontsize{11}{13.86}\selectfont 9 shared sides};
\node[anchor=north west,inner sep=0pt,outer sep=0pt,align=left] at (15,118) {\fontsize{11}{13.86}\selectfont Boundary: \rule{14mm}{.3pt}};
\begin{scope}[shift={(114,77)}]
\draw[step=6.2mm,black!27,line width=.3pt] (0,0) grid (74.4,31.0);
\end{scope}
\draw[black!27,line width=.3pt] (114,77) rectangle (188.4,108.0);
\node[anchor=north west,inner sep=0pt,outer sep=0pt,align=left] at (114,111) {\fontsize{11}{13.86}\selectfont 10 shared sides};
\node[anchor=north west,inner sep=0pt,outer sep=0pt,align=left] at (114,118) {\fontsize{11}{13.86}\selectfont Boundary: \rule{14mm}{.3pt}};
\begin{scope}[shift={(15,132)}]
\draw[step=6.2mm,black!27,line width=.3pt] (0,0) grid (74.4,31.0);
\end{scope}
\draw[black!27,line width=.3pt] (15,132) rectangle (89.4,163.0);
\node[anchor=north west,inner sep=0pt,outer sep=0pt,align=left] at (15,166) {\fontsize{11}{13.86}\selectfont 11 shared sides};
\node[anchor=north west,inner sep=0pt,outer sep=0pt,align=left] at (15,173) {\fontsize{11}{13.86}\selectfont Boundary: \rule{14mm}{.3pt}};
\begin{scope}[shift={(114,132)}]
\draw[step=6.2mm,black!27,line width=.3pt] (0,0) grid (74.4,31.0);
\end{scope}
\draw[black!27,line width=.3pt] (114,132) rectangle (188.4,163.0);
\node[anchor=north west,inner sep=0pt,outer sep=0pt,align=left] at (114,166) {\fontsize{11}{13.86}\selectfont 12 shared sides};
\node[anchor=north west,inner sep=0pt,outer sep=0pt,align=left] at (114,173) {\fontsize{11}{13.86}\selectfont Boundary: \rule{14mm}{.3pt}};
\begin{scope}[shift={(15,187)}]
\draw[step=6.2mm,black!27,line width=.3pt] (0,0) grid (74.4,31.0);
\end{scope}
\draw[black!27,line width=.3pt] (15,187) rectangle (89.4,218.0);
\node[anchor=north west,inner sep=0pt,outer sep=0pt,align=left] at (15,221) {\fontsize{11}{13.86}\selectfont 13 shared sides};
\node[anchor=north west,inner sep=0pt,outer sep=0pt,align=left] at (15,228) {\fontsize{11}{13.86}\selectfont Boundary: \rule{14mm}{.3pt}};
\draw[black!30,line width=.35pt] (114,196) -- ++(80,0);
\draw[black!30,line width=.35pt] (114,211) -- ++(80,0);
\draw[black!30,line width=.35pt] (114,226) -- ++(80,0);
\draw[black!30,line width=.35pt] (15,246) -- ++(179,0);
\draw[black!30,line width=.35pt] (15,259) -- ++(179,0);
\node[anchor=north west,inner sep=0pt,outer sep=0pt,align=left] at (15,268) {\fontsize{9}{11.34}\selectfont Bellingham Math Circle / Week 26 / W26-23-v2};
\node[anchor=north east,inner sep=0pt,outer sep=0pt,align=left] at (199,268) {\fontsize{9}{11.34}\selectfont 3};
\end{scope}
\end{tikzpicture}
\newpage
\null
\begin{tikzpicture}[remember picture,overlay,x=1mm,y=-1mm]
\begin{scope}[shift={(current page.north west)}]
\node[anchor=north west,inner sep=0pt,outer sep=0pt,align=left,text width=187mm] at (15,12) {\fontsize{11}{13.86}\selectfont Week 26 / Same area, different boundaries / Grades 2--3};
\node[anchor=north west,inner sep=0pt,outer sep=0pt,align=left,text width=184mm] at (15,28) {\fontsize{14}{17.64}\selectfont \textbf{Problem 4:} Find the longest possible boundary with 4, 7, 10, and 12 tiles. Explain why no shape using each number of tiles can have a longer boundary.};
\begin{scope}[shift={(15,67)}]
\draw[step=5.5mm,black!27,line width=.3pt] (0,0) grid (77.0,27.5);
\end{scope}
\draw[black!27,line width=.3pt] (15,67) rectangle (92.0,94.5);
\node[anchor=north west,inner sep=0pt,outer sep=0pt,align=left] at (15,98) {\fontsize{11}{13.86}\selectfont 4 tiles; boundary: \rule{12mm}{.3pt}};
\begin{scope}[shift={(114,67)}]
\draw[step=5.5mm,black!27,line width=.3pt] (0,0) grid (77.0,27.5);
\end{scope}
\draw[black!27,line width=.3pt] (114,67) rectangle (191.0,94.5);
\node[anchor=north west,inner sep=0pt,outer sep=0pt,align=left] at (114,98) {\fontsize{11}{13.86}\selectfont 7 tiles; boundary: \rule{12mm}{.3pt}};
\begin{scope}[shift={(15,133)}]
\draw[step=5.5mm,black!27,line width=.3pt] (0,0) grid (77.0,27.5);
\end{scope}
\draw[black!27,line width=.3pt] (15,133) rectangle (92.0,160.5);
\node[anchor=north west,inner sep=0pt,outer sep=0pt,align=left] at (15,164) {\fontsize{11}{13.86}\selectfont 10 tiles; boundary: \rule{12mm}{.3pt}};
\begin{scope}[shift={(114,133)}]
\draw[step=5.5mm,black!27,line width=.3pt] (0,0) grid (77.0,27.5);
\end{scope}
\draw[black!27,line width=.3pt] (114,133) rectangle (191.0,160.5);
\node[anchor=north west,inner sep=0pt,outer sep=0pt,align=left] at (114,164) {\fontsize{11}{13.86}\selectfont 12 tiles; boundary: \rule{12mm}{.3pt}};
\draw[black!30,line width=.35pt] (15,195) -- ++(179,0);
\draw[black!30,line width=.35pt] (15,215) -- ++(179,0);
\draw[black!30,line width=.35pt] (15,235) -- ++(179,0);
\draw[black!30,line width=.35pt] (15,255) -- ++(179,0);
\node[anchor=north west,inner sep=0pt,outer sep=0pt,align=left] at (15,268) {\fontsize{9}{11.34}\selectfont Bellingham Math Circle / Week 26 / W26-23-v2};
\node[anchor=north east,inner sep=0pt,outer sep=0pt,align=left] at (199,268) {\fontsize{9}{11.34}\selectfont 4};
\end{scope}
\end{tikzpicture}
\newpage
\null
\begin{tikzpicture}[remember picture,overlay,x=1mm,y=-1mm]
\begin{scope}[shift={(current page.north west)}]
\node[anchor=north west,inner sep=0pt,outer sep=0pt,align=left,text width=187mm] at (15,12) {\fontsize{11}{13.86}\selectfont Week 26 / Same area, different boundaries / Grades 2--3};
\node[anchor=north west,inner sep=0pt,outer sep=0pt,align=left,text width=184mm] at (15,28) {\fontsize{14}{17.64}\selectfont \textbf{Problem 5:} Which of these eight-tile shapes have the longest possible boundary? Explain why each chosen shape reaches that length. Draw two more such shapes, with no straight row or column of four tiles.};
\filldraw[fill=black!13,draw=black,line width=.75pt] (15.0,86.0) rectangle ++(7.5,7.5);
\filldraw[fill=black!13,draw=black,line width=.75pt] (22.5,86.0) rectangle ++(7.5,7.5);
\filldraw[fill=black!13,draw=black,line width=.75pt] (30.0,86.0) rectangle ++(7.5,7.5);
\filldraw[fill=black!13,draw=black,line width=.75pt] (37.5,86.0) rectangle ++(7.5,7.5);
\filldraw[fill=black!13,draw=black,line width=.75pt] (45.0,86.0) rectangle ++(7.5,7.5);
\filldraw[fill=black!13,draw=black,line width=.75pt] (52.5,86.0) rectangle ++(7.5,7.5);
\filldraw[fill=black!13,draw=black,line width=.75pt] (60.0,86.0) rectangle ++(7.5,7.5);
\filldraw[fill=black!13,draw=black,line width=.75pt] (67.5,86.0) rectangle ++(7.5,7.5);
\filldraw[fill=black!13,draw=black,line width=.75pt] (114.0,85.5) rectangle ++(7.5,7.5);
\filldraw[fill=black!13,draw=black,line width=.75pt] (121.5,85.5) rectangle ++(7.5,7.5);
\filldraw[fill=black!13,draw=black,line width=.75pt] (129.0,85.5) rectangle ++(7.5,7.5);
\filldraw[fill=black!13,draw=black,line width=.75pt] (136.5,85.5) rectangle ++(7.5,7.5);
\filldraw[fill=black!13,draw=black,line width=.75pt] (144.0,85.5) rectangle ++(7.5,7.5);
\filldraw[fill=black!13,draw=black,line width=.75pt] (114.0,78.0) rectangle ++(7.5,7.5);
\filldraw[fill=black!13,draw=black,line width=.75pt] (129.0,93.0) rectangle ++(7.5,7.5);
\filldraw[fill=black!13,draw=black,line width=.75pt] (144.0,78.0) rectangle ++(7.5,7.5);
\filldraw[fill=black!13,draw=black,line width=.75pt] (29.0,132.0) rectangle ++(7.5,7.5);
\filldraw[fill=black!13,draw=black,line width=.75pt] (29.0,139.5) rectangle ++(7.5,7.5);
\filldraw[fill=black!13,draw=black,line width=.75pt] (36.5,132.0) rectangle ++(7.5,7.5);
\filldraw[fill=black!13,draw=black,line width=.75pt] (36.5,139.5) rectangle ++(7.5,7.5);
\filldraw[fill=black!13,draw=black,line width=.75pt] (44.0,132.0) rectangle ++(7.5,7.5);
\filldraw[fill=black!13,draw=black,line width=.75pt] (44.0,139.5) rectangle ++(7.5,7.5);
\filldraw[fill=black!13,draw=black,line width=.75pt] (51.5,132.0) rectangle ++(7.5,7.5);
\filldraw[fill=black!13,draw=black,line width=.75pt] (51.5,139.5) rectangle ++(7.5,7.5);
\filldraw[fill=black!13,draw=black,line width=.75pt] (134.0,127.0) rectangle ++(7.5,7.5);
\filldraw[fill=black!13,draw=black,line width=.75pt] (134.0,134.5) rectangle ++(7.5,7.5);
\filldraw[fill=black!13,draw=black,line width=.75pt] (134.0,142.0) rectangle ++(7.5,7.5);
\filldraw[fill=black!13,draw=black,line width=.75pt] (141.5,127.0) rectangle ++(7.5,7.5);
\filldraw[fill=black!13,draw=black,line width=.75pt] (141.5,142.0) rectangle ++(7.5,7.5);
\filldraw[fill=black!13,draw=black,line width=.75pt] (149.0,127.0) rectangle ++(7.5,7.5);
\filldraw[fill=black!13,draw=black,line width=.75pt] (149.0,134.5) rectangle ++(7.5,7.5);
\filldraw[fill=black!13,draw=black,line width=.75pt] (149.0,142.0) rectangle ++(7.5,7.5);
\begin{scope}[shift={(15,188)}]
\draw[step=7.5mm,black!27,line width=.3pt] (0,0) grid (75.0,37.5);
\end{scope}
\draw[black!27,line width=.3pt] (15,188) rectangle (90.0,225.5);
\begin{scope}[shift={(114,188)}]
\draw[step=7.5mm,black!27,line width=.3pt] (0,0) grid (75.0,37.5);
\end{scope}
\draw[black!27,line width=.3pt] (114,188) rectangle (189.0,225.5);
\draw[black!30,line width=.35pt] (15,241) -- ++(179,0);
\draw[black!30,line width=.35pt] (15,257) -- ++(179,0);
\node[anchor=north west,inner sep=0pt,outer sep=0pt,align=left] at (15,268) {\fontsize{9}{11.34}\selectfont Bellingham Math Circle / Week 26 / W26-23-v2};
\node[anchor=north east,inner sep=0pt,outer sep=0pt,align=left] at (199,268) {\fontsize{9}{11.34}\selectfont 5};
\end{scope}
\end{tikzpicture}
\newpage
\null
\begin{tikzpicture}[remember picture,overlay,x=1mm,y=-1mm]
\begin{scope}[shift={(current page.north west)}]
\node[anchor=north west,inner sep=0pt,outer sep=0pt,align=left,text width=187mm] at (15,12) {\fontsize{11}{13.86}\selectfont Week 26 / Same area, different boundaries / Grades 2--3};
\node[anchor=north west,inner sep=0pt,outer sep=0pt,align=left,text width=184mm] at (15,28) {\fontsize{14}{17.64}\selectfont \textbf{Problem 6:} Add one tile to a shape so that its boundary changes by each amount shown. For each amount, use as few starting tiles as possible. Explain why fewer tiles cannot work.};
\begin{scope}[shift={(22,82)}]
\draw[step=8mm,black!27,line width=.3pt] (0,0) grid (64,40);
\end{scope}
\draw[black!27,line width=.3pt] (22,82) rectangle (86,122);
\node[anchor=north west,inner sep=0pt,outer sep=0pt,align=left] at (22,126) {\fontsize{12}{15.12}\selectfont Boundary change: +2};
\node[anchor=north west,inner sep=0pt,outer sep=0pt,align=left] at (22,135) {\fontsize{12}{15.12}\selectfont Starting tiles: \rule{10mm}{.3pt}};
\begin{scope}[shift={(121,82)}]
\draw[step=8mm,black!27,line width=.3pt] (0,0) grid (64,40);
\end{scope}
\draw[black!27,line width=.3pt] (121,82) rectangle (185,122);
\node[anchor=north west,inner sep=0pt,outer sep=0pt,align=left] at (121,126) {\fontsize{12}{15.12}\selectfont Boundary change: 0};
\node[anchor=north west,inner sep=0pt,outer sep=0pt,align=left] at (121,135) {\fontsize{12}{15.12}\selectfont Starting tiles: \rule{10mm}{.3pt}};
\begin{scope}[shift={(22,170)}]
\draw[step=8mm,black!27,line width=.3pt] (0,0) grid (64,40);
\end{scope}
\draw[black!27,line width=.3pt] (22,170) rectangle (86,210);
\node[anchor=north west,inner sep=0pt,outer sep=0pt,align=left] at (22,214) {\fontsize{12}{15.12}\selectfont Boundary change: $-2$};
\node[anchor=north west,inner sep=0pt,outer sep=0pt,align=left] at (22,223) {\fontsize{12}{15.12}\selectfont Starting tiles: \rule{10mm}{.3pt}};
\begin{scope}[shift={(121,170)}]
\draw[step=8mm,black!27,line width=.3pt] (0,0) grid (64,40);
\end{scope}
\draw[black!27,line width=.3pt] (121,170) rectangle (185,210);
\node[anchor=north west,inner sep=0pt,outer sep=0pt,align=left] at (121,214) {\fontsize{12}{15.12}\selectfont Boundary change: $-4$};
\node[anchor=north west,inner sep=0pt,outer sep=0pt,align=left] at (121,223) {\fontsize{12}{15.12}\selectfont Starting tiles: \rule{10mm}{.3pt}};
\draw[black!30,line width=.35pt] (15,240) -- ++(179,0);
\draw[black!30,line width=.35pt] (15,256) -- ++(179,0);
\node[anchor=north west,inner sep=0pt,outer sep=0pt,align=left] at (15,268) {\fontsize{9}{11.34}\selectfont Bellingham Math Circle / Week 26 / W26-23-v2};
\node[anchor=north east,inner sep=0pt,outer sep=0pt,align=left] at (199,268) {\fontsize{9}{11.34}\selectfont 6};
\end{scope}
\end{tikzpicture}
\end{document}
